\centerline{\bf \Large  8 класс}

\problem{\textbf{1}} Пюрвя назвал натуральное число \textit{УДЕкрасивым}, если сумма цифр числа равна 20, при этом никакая цифра не встречается дважды. Сколько существует \textit{УДЕкрасивых} чисел, меньших 2025?

\textbf{Решение:} Выпишем все возможные варианты УДЕкрасивых чисел:

389, 398

479, 497

569, 578, 587, 596

659, 695

749, 758, 785, 794

839, 857, 875, 893

938, 947, 956, 965, 974, 983

1289, 1298

1379, 1397

1469, 1478, 1487, 1496

1568, 1586

1649, 1658, 1685, 1694

1739, 1748, 1784, 1793

1829, 1847, 1856, 1865, 1874, 1892

1928, 1937, 1946, 1964, 1973, 1982

Всего 54 чисел.

\textbf{Критерии.} \\
\textit{7 баллов} $-$ 53-54 правильных чисел\\
\textit{6 баллов} $-$ 44-52 правильных чисел\\
\textit{5 баллов} $-$ 36-43 правильных чисел\\
\textit{4 балла} $-$ 29-35 правильных чисел\\
\textit{3 балла} $-$ 22-28 правильных чисел\\
\textit{2 балла} $-$ 15-21 правильных чисел\\
\textit{1 балл} $-$ 8-14 правильных чисел\\
\textit{0 баллов} $-$ 0-7 правильных чисел\\
\textit{0-7 баллов} $-$ Комбинаторный способ\\

\problem{\textbf{2}} Решите задачу:

а) Батыр и Очир ходили вокруг школы, оба в одном направлении и каждый с постоянной скоростью. Батыр проходил один круг за 3 минуты, а Очир, который был медленнее Батыра, - за некоторое целое число минут. При этом время между встречами тоже равнялось некоторому целому числу минут, причём оно было больше 7. За какое время Очир проходил полный круг?


б) Составьте и решите обратную задачу по схеме: ?; 4; 7.

\textbf{Решение А:} Обозначим через $n$ время в минутах, за которое проходит круг Очир. Тогда $n>3, n \in \mathbb{N}$. Скорость, с которой более быстрый мальчик догоняет медленного равна

$$
\frac{1}{3}-\frac{1}{n}=\frac{n-3}{3 n}\left(\frac{\kappa p y r}{min}\right),
$$


поэтому время между встречами составляет

$$
\frac{3 n}{n-3}=\frac{3(n-3)+9}{n-3}=3+\frac{9}{n-3}(\text { мин }),
$$


а поскольку это целое число, то $(n-3)$ делит 9. Натуральные делители 9 - числа $1 , 3$ и 9 . Тогда $n=12$ или $n=6$ или $n=4$, что соответствует времени между встречами 4 минут, 6 минут и 12 минут. Поскольку время между встречами больше 7 минут, заданному условию удовлетворяет только $n=4$.



\textbf{Примерное условие Б:} Батыр и Очир ходили вокруг школы, оба в одном направлении и каждый с постоянной скоростью. Очир проходил один круг за 4 минуты, а Батыр, который был быстрее Очира, - за некоторое целое число минут. При этом время между встречами тоже равнялось некоторому целому числу минут, причём оно было больше 7. За какое время Очир проходил полный круг?

\textbf{Решение Б:} Обозначим через $n$ время в минутах, за которое проходит круг Батыр. Тогда $n<4, n \in \mathbb{N}$. Скорость, с которой более быстрый мальчик догоняет медленного равна

$$
\frac{1}{n}-\frac{1}{4}=\frac{4-n}{4 n}\left(\frac{\kappa p y r}{min}\right),
$$


поэтому время между встречами составляет

$$
\frac{4 n}{4-n}=\frac{16-(4-n)4}{4-n}=\frac{16}{4-n}-4(\text { мин }),
$$


а поскольку это целое число, то $(4-n)$ делит 16 и $\dfrac{16}{4-n}-4>0$. Тогда $n=2$ или $n=3$, что соответствует времени между встречами 4 минут и 12 минут. Поскольку время между встречами больше 7 минут, заданному условию удовлетворяет только $n=3$.

\textbf{Критерии.} \\
\textit{0-4 баллов} $-$ Решение пункта А\\
\textit{0-1 балла} $-$ Составление условия пункта Б \\
\textit{0-2 баллов} $-$ Решение пункта Б\\
\textit{0 баллов} $-$ Любой ответ без объяснений\\


\problem{\textbf{3}} В стране УДЕ мальчики днём всегда говорят правду, а ночью всегда лгут, а девочки — наоборот. Днём каждый из них сказал: “У меня знакомых мальчиков на 1 больше, чем знакомых девочек”, а ночью — “Среди незнакомых мне жителей страны мальчиков на 1 больше, чем девочек”. Могло ли в стране УДЕ быть ровно 2025 жителей?

\textbf{Решение:} Для решения задачи полезно знать следующий факт, который мы для удобства сформулируем в терминах знакомства. Не существует компании, состоящей из нечётного числа людей, в которой каждый человек имеет нечётное количество знакомых (среди людей из этой компании).

Доказательство. Пусть все знакомые пожмут друг другу руки. Так как в каждом рукопожатии участвуют две руки, значит, суммарное количество «протянутых рук» четно. С другой стороны, каждый человек протягивал руку нечётное число раз, Чтобы суммарное число «протянутых рук» в этой ситуации получилось чётным, число людей должно быть чётным.

Каждый мальчик в днём говорил правду, значит, у него нечётное число знакомых (поскольку среди них мальчиков на 1 больше, чем девочек). А каждая девочка сказала правду ночью, и по той же причине у неё нечётное число незнакомых.

Предположим, что в стране 2025 жителей. Тогда у каждого человека сумма количества знакомых и незнакомых равна 2024 (чётное число). Следовательно, у каждого человека количество знакомых имеет ту же чётность, что и количество незнакомых.

Таким образом, у каждой девочки нечётное количество незнакомых и, следовательно, нечётное количество знакомых. В итоге получаем, что каждый из 2025 жителей, имеет нечётное число знакомых, чего быть не может.

\textbf{Критерии.} \\
\textit{0-5 баллов} $-$ Продвижение через идеи чётности\\
\textit{1 балл} $-$ Лемма о рукопожатиях без доказательства\\
\textit{1 балл} $-$ Доказательство леммы о рукопожатиях\\
\textit{0 баллов} $-$ Любой ответ без объяснений\\

\problem{\textbf{4}} Существует ли треугольник, в котором одна сторона равна какой-то из его высот, другая - какой-то из биссектрис, а третья - какой-то из медиан?

\textbf{Решение:} Нет, так как наибольшая сторона треугольника длиннее любой из его высот, медиан или биссектрис. Действительно, любой отрезок, соединяющий вершину треугольника с точкой на противоположной стороне, короче, по крайней мере, одной из двух других сторон. Поэтому любая медиана или биссектриса короче хотя бы одной из сторон и, тем самым, короче наибольшей стороны. Тем более, это верно для высот.

\textbf{Критерии.} \\
\textit{7 баллов} $-$ Полностью верное решение\\
\textit{0-6 баллов} $-$ Все остальные продвижения\\
\textit{0 баллов} $-$ Любой ответ без объяснений\\

\problem{\textbf{5}} 
Назовём число $n$ - \textit{УДЕинтересным}, если существует точный квадрат, сумма цифр которого равна $n$. Найдите количество \textit{УДЕинтересных} чисел от $1$ до $2025$.

\textbf{Решение:} Квадрат любого числа либо делится на 9 , либо даёт при делении на 9 один из остатков 1,4 и 7. Поскольку сумма цифр любого натурального числа даёт при делении на 9 тот же остаток, что и само число, то сумма цифр квадрата может быть равна $9 n, 9 n+1,9 n+$ +4 или $9 n+7$, где $n-$ целое неотрицательное число. Докажем, что все такие числа могут быть суммами:

$$
\begin{aligned}
9 n & \text { - сумма цифр числа }\left(10^n-1\right)^2=10^{2 n}-2 \cdot 10^n+1= \\
& =\underbrace{9 \ldots 9}_{n-1} \underbrace{0 \ldots 0}_{n-1} 1 .
\end{aligned}
$$


$$
\text { очевиден: сумма цифр числа } 1^2 \text { равна } 1 .
$$

$9 n+4$, где $n>0$, - сумма цифр числа $\left(10^n-3\right)^2=$ $=10^{2 n}-6 \cdot 10^n+9=\underbrace{9 \ldots 9}_{n-1}\ \underbrace{0 \ldots 0}_{n-1} 9$. Случай $n=0$ опять очевиден: сумма цифр числа $2^2$ равна 4.

$$
\begin{aligned}
& 9 n+7-\text { сумма цифр числа }\left(10^{n+1}-5\right)^2=10^{2 n+2}- \\
& \quad-10^{n+2}+25=\underbrace{9 \ldots 9}_n \underbrace{0 \ldots 0}_n 25 .
\end{aligned}
$$

В таком случае от 1 до 2025 ровно $\dfrac{4}{5}\cdot2025 = 900$ УДЕинтересных чисел.

\textbf{Критерии.} \\
\textit{7 баллов} $-$ Полностью верное решение\\
\textit{0-6 баллов} $-$ Все остальные продвижения\\
\textit{0 баллов} $-$ Любой ответ без объяснений\\